\documentclass[11pt,a4paper]{article}

\usepackage[utf8]{inputenc}
\usepackage[T1]{fontenc}
\usepackage{amsmath,amssymb,amsthm}
\usepackage{booktabs}
\usepackage{enumitem}
\usepackage{microtype}
\usepackage[margin=2.5cm]{geometry}
\usepackage{hyperref}
\hypersetup{colorlinks=true,linkcolor=blue,citecolor=blue}

\newcommand{\G}{\mathcal{G}}
\newcommand{\range}{\mathrm{range}}
\newcommand{\Types}{\mathrm{Types}}

\title{Beyond Traditional KGQA:\\What LLMs Are For, and Why the\\
Field Is Genuinely Hard}
\author{Bernard Katamanso}
\date{\today}

\begin{document}
\maketitle

%======================================================================
\section{What Do We Actually Want LLMs For?}
%======================================================================

The common failure mode is to treat the LLM as a \emph{knowledge
store}.  It is not: parametric memory is frozen at training time,
unverifiable, and hallucinates precisely where facts are
domain-specific.  The useful reframing is a \emph{division of labor}
between a linguistic system and a structured system.

\begin{center}
\begin{tabular}{@{}lll@{}}
\toprule
  Role & Performed by & Why the LLM \\
\midrule
  \textbf{Comprehend} & LLM & Turn noisy natural language into intent,
    entity mentions, and constraints \\
  \textbf{Compose} & LLM + symbolic layer & Plan/execute multi-hop reasoning:
    the LLM \emph{proposes}, structure \emph{disposes} \\
  \textbf{Communicate} & LLM & Produce the answer plus explanation in
    natural language \\
\bottomrule
\end{tabular}
\end{center}

The knowledge graph is the ground truth; the LLM is the \emph{linguistic
front-end} and \emph{reasoning orchestrator}---never the memory.  This
is the neuro-symbolic split, and DCA-Trie instantiates it: the LLM
generates, the TypeOracle and trie decide what is legal.

The LLM earns its place exactly when the question is \textbf{not
answerable by lookup alone}: when it needs composition, path-finding,
or a natural-language interface to structured data.  Users of a
substation or hospital system will not write SPARQL; they will say
``which feeder is feeding the faulted bus?''  That utterance requires
comprehension, decomposition, and path-finding over a structured model.

%======================================================================
\section{RAG vs.\ Agentic Retrieval vs.\ Constrained Decoding}
%======================================================================

These are three points on the \textbf{fluency--fidelity tradeoff}:
how tightly the model's output is bound to external structure.

\subsection{Retrieval-Augmented Generation (RAG)}

\begin{itemize}
  \item \emph{Context} is grounded: retrieve relevant passages/subgraph,
    then generate conditioned on them.
  \item The \emph{generation} remains free: the model can still ignore
    the retrieved context and answer from parametric memory.
  \item One-shot, no feedback loop.
  \item Retrieval quality is the bottleneck: bad retrieval $\Rightarrow$
    bad answer, regardless of generation quality.
\end{itemize}

\subsection{Agentic Retrieval}

\begin{itemize}
  \item Iterative, tool-using, reflective: the model decides what to
    fetch, observes the results, and re-plans.
  \item Buys \textbf{backtracking and recovery}---the structural
    capability that a single-pass decoder lacks.
  \item Examples: KBQA-o1 (MCTS), GraphRAG-style agent loops.
  \item Cost: reasoning loops can explode; tool selection is still
    model-decided, so the agent can chase the wrong subgraph.
\end{itemize}

\subsection{Constrained Decoding}

\begin{itemize}
  \item The tightest coupling: the graph constrains the
    \emph{generation process itself} at the token level.
  \item RAG \emph{suggests}; constrained decoding \emph{forbids}.
    It is impossible for the model to emit a token that leads off the
    graph.
  \item The strongest fidelity guarantee---at the cost of
    \textbf{dead ends}: over-constraint can leave no valid completion.
\end{itemize}

\begin{center}
\begin{tabular}{@{}lccc@{}}
\toprule
  & RAG & Agentic & Constrained \\
\midrule
  Grounding point & Context & Process & Tokens \\
  Feedback loop    & No      & Yes      & No \\
  Backtracking     & No      & Yes      & No \\
  Can hallucinate off-graph & Yes & Yes & \textbf{No} \\
  Risk             & Stale/irrelevant context & Loop explosion & Dead ends \\
\bottomrule
\end{tabular}
\end{center}

The interesting design space is \emph{combining} them: agentic planning
on the outside, constrained decoding on the inside.

%======================================================================
\section{You Can Define Your Own Schema---and That Is the Payoff}
%======================================================================

This is the key reason the TypeOracle generalizes beyond Freebase.
Freebase is a \emph{general, shallow} ontology, which is exactly why
our gates are permissive (13\% path reduction).  A \textbf{domain
schema is narrow, deep, and precise}, so semantic gates become much
stronger.

\subsection{Substation / Electrical Grid (IEC 61850, IEC 61970)}

\begin{itemize}
  \item \textbf{Entities:} transformers, breakers, buses, relays,
    protection zones, feeders.
  \item \textbf{Relations:} \texttt{feeds}, \texttt{protects},
    \texttt{measures}, \texttt{connected\_to}.
  \item Example gate: $\range(\texttt{protects}) = \text{protection
    devices} \to \text{protection zones}$.  The path
    \texttt{transformer $\to$ protects $\to$ bus} is rejected because a
    transformer is not a protection device.  In a safety-critical
    domain, that rejection is the whole point.
\end{itemize}

\subsection{Health (SNOMED CT, RxNorm, ICD)}

\begin{itemize}
  \item \textbf{Entities:} patients, conditions, drugs, genes,
    proteins, trials.
  \item \textbf{Relations:} \texttt{treats}, \texttt{contraindicates},
    \texttt{associated\_with}.
  \item Example gate: a \texttt{contraindicates} gate prevents pairing a
    drug with an incompatible condition.
\end{itemize}

\subsection{The Living Graph}

A substation's SCADA telemetry is an event stream that \emph{mutates}
the graph: breakers open, loads change, zones reconfigure.  The
constrained-decoding approach survives this because the trie is rebuilt
from the \textbf{current snapshot} for each query.  The questions become
temporal---``what is the state of the network \emph{now}?''---which is
temporal KGQA, a regime where a frozen parametric model fundamentally
cannot compete.

That the schema is \emph{yours} is also the point: Freebase constrains
weakly because it is generic; your domain ontology constrains strongly
because you control it.

%======================================================================
\section{Why the Field Is Genuinely Hard}
%======================================================================

Seven fundamental barriers---not engineering issues.

\subsection{1. The representation gap}
Natural language is linear, ambiguous, and under-specified; a graph is
exact and combinatorial.  Mapping ``the big fuse box thing'' to
\texttt{transformer \#7} is entity linking, and it is genuinely
unsolved for messy input.

\subsection{2. Combinatorial reasoning}
Multi-hop reasoning is exponential: $O(d^L)$ paths.  Every constraint
that makes this tractable trades away recall.  This is an
information-theoretic tension, not a tuning problem.

\subsection{3. Open-world incompleteness}
KGs miss edges, miss types, and contain aliases and errors.  The
oracle's conservatism (``allow when unknown'') is a double-edged
sword: it never wrongly filters, but it barely filters at all.

\subsection{4. Ontology quality is the ceiling}
Semantic gates are only as strong as the type system.  Coarse
ontology $\Rightarrow$ weak pruning (our 13\% on Freebase).  This is a
\emph{data} problem, not an algorithm problem.

\subsection{5. Hallucination is the model's inductive bias}
LLMs are trained to be fluent, and fluency means ``plausible
continuation''---not ``faithful to the graph.''  Constraining
generation fights the training objective itself.  This is why
constrained decoding is genuinely harder than it sounds.

\subsection{6. Living knowledge}
Static tries and embeddings go stale.  Reasoning over changing graphs
(temporal KGQA) is largely unsolved.

\subsection{7. Trust, not accuracy}
89\% Hits@1 is a publication result; it is \emph{unacceptable} for a
hospital or a grid operator.  Those domains demand provable, citable,
explainable answers---and the ability to say ``I do not know'' rather
than guess.  The entire fidelity apparatus exists to serve
\emph{verifiability}, and current evaluation metrics do not even
measure it.

%======================================================================
\section{Synthesis}
%======================================================================

The field is hard because it asks a \textbf{probabilistic, ungrounded
system} to behave like a \textbf{deterministic, grounded query
engine}---while preserving the fluency that makes the probabilistic
system useful in the first place.

Every method is a point on the fluency--fidelity tradeoff:

\[
  \underbrace{\text{Pure LLM}}_{\max \text{ fluency},\; \min \text{ fidelity}}
  \;\longrightarrow\;
  \underbrace{\text{RAG}}_{\text{context grounded}}
  \;\longrightarrow\;
  \underbrace{\text{Agentic}}_{\text{process grounded}}
  \;\longrightarrow\;
  \underbrace{\text{Constrained decoding}}_{\min \text{ fluency},\; \max \text{ fidelity}}
\]

DCA-Trie's thesis is that one can have \emph{a lot} of both by
constraining at the token level with a \emph{semantic} gate---not just
a topological one---and especially by controlling the schema.  The
concrete move from 13\% Freebase pruning to a 60\%+ reduction on a
domain ontology is the empirical path forward: the ceiling on
pruning is set by the ontology, and the ontology is now ours to build.

\end{document}
